\section{Agent 在实践中的应用领域}

文章附录详细讨论了两个最有前景的 Agent 应用领域，它们共享以下特征：需要对话与行动的结合、有明确的成功标准、能形成反馈循环、并整合有意义的人类监督。

\subsection{客户支持（Customer Support）}

客户支持是 Agent 的天然契合场景：

\begin{itemize}
    \item 支持交互本身就是对话流，同时需要访问外部信息和执行操作。
    \item 可以集成工具来拉取客户数据、订单历史、知识库文章。
    \item 退款、更新工单等操作可以程序化处理。
    \item 成功可以通过用户确认的解决方案来明确衡量。
\end{itemize}

\begin{knowledgebox}{商业模式验证}
多家公司已通过\textbf{基于成功解决方案收费}（usage-based pricing）的商业模式验证了这一方向——只有当 Agent 成功解决用户问题时才收费。这种商业模式本身就是对 Agent 能力的最强信号。
\end{knowledgebox}

\subsection{编程 Agent（Coding Agents）}

软件开发领域展现了 LLM Agent 的巨大潜力，能力从代码补全进化到自主问题解决。编程 Agent 之所以特别有效，原因在于：

\begin{itemize}
    \item 代码解决方案可以通过\textbf{自动化测试}验证。
    \item Agent 可以用测试结果作为反馈进行迭代。
    \item 问题空间定义明确、结构化。
    \item 输出质量可以客观衡量。
\end{itemize}

\begin{warningbox}{自动化测试不等于完全验证}
虽然自动化测试能验证代码的功能正确性，但\textbf{人类审查}仍然不可或缺——需要确保解决方案符合更广泛的系统需求、代码风格规范、性能要求、安全标准等。测试通过只是最低门槛，不是质量保证的全部。
\end{warningbox}

\subsection{本章小结}
客户支持和编程是目前最成熟的 Agent 应用领域。前者契合于对话式交互+操作的模式，后者受益于代码的可测试性和问题空间的结构化。两者都需要人类监督来保证更高层次的质量。


